
\subsubsection{3.1 Problem Setup}

Consider a classification task with training data $\{(x_i, y_i, g_i)\}_{i=1}^{N}$ where $x_i$ is the input, $y_i \in \{0, 1\}$ is the class label, and $g_i \in \{1, 2, 3, 4\}$ encodes the combination of class and spurious attribute. In Waterbirds, groups are defined by (bird type $\times$ background): majority groups have aligned spurious features (waterbirds on water, landbirds on land), while minority groups have conflicting features (waterbirds on land, landbirds on water).

Under standard ERM training, the model minimizes aggregate loss without access to group labels:

$L_{\mathrm{ERM}} = (1/N) \sum_i \ell(f_\theta(x_i), y_i)$

The goal is to determine whether the per-sample loss trajectories $\{\ell_i^{(t)}\}_{t=1}^{T}$ during ERM training encode information distinguishing minority from majority samples.

\subsubsection{3.2 Per-Sample Loss Tracking}

Per-sample cross-entropy loss is tracked across training epochs. To reduce stochastic noise, evaluation passes after each epoch are run with data augmentation disabled and batch ordering fixed. Individual sample losses are computed using reduction='none' in the loss function. Trajectory features are extracted from epochs 1-5.

\subsubsection{3.3 Trajectory Feature Extraction}

From the loss matrix $L \in \mathbb{R}^{N \times T}$ where $L_{i,t}$ is sample $i$'s loss at epoch $t$, four features are extracted per sample:

\begin{itemize}
\item \textbf{Initial Loss ($L_{1})$:} The loss value at epoch 1, capturing immediate conflict between the sample and early model representations.
\item \textbf{Slope:} The linear trend of loss decrease across epochs, capturing learning speed.
\item \textbf{Variance:} The variability of losses across epochs, capturing trajectory instability.
\item \textbf{Convergence Time:} The first epoch where normalized loss drops below a threshold $\tau$ = 0.1.
\end{itemize}

\subsubsection{3.4 Minority Prediction Evaluation}

Trajectory features are evaluated for their ability to predict minority group membership using logistic regression with 5-fold stratified cross-validation. AUROC is the primary metric, with a pre-specified success criterion of AUROC > 0.75.

\subsubsection{3.5 Spurious-Specificity Test}

To test whether trajectory divergence reflects spurious correlation specifically rather than generic sample difficulty, a controlled experiment compares three training regimes:

\begin{itemize}
\item \textbf{ERM (baseline):} Standard empirical risk minimization.
\item \textbf{GroupDRO:} Minimizes worst-group loss, specifically targeting spurious correlation reliance.
\item \textbf{Variance-matched random reweighting:} A control condition matching GroupDRO's gradient variance without targeting spurious correlations.
\end{itemize}

If trajectory divergence is spurious-specific, AUROC should decrease substantially under GroupDRO ($\Delta$ > 0.10) but remain stable under random reweighting ($\Delta$ < 0.05).

\subsubsection{3.6 Curvature Timing Analysis}

A secondary hypothesis proposed that minority samples exhibit delayed curvature stabilization---the epoch at which the second derivative of the loss curve transitions from negative to near-zero. This was tested by computing curvature via central differences on smoothed loss curves (Gaussian $\sigma$ = 1.0) and detecting the sign-flip epoch where curvature exceeds -0.002 for two consecutive epochs. The hypothesis predicted a timing gap of $\geq 3$ epochs between minority and majority groups in $\geq 70$\% of seeds.
